\MathBookSourcePage{177}
\subsection{第 $k$ 层迭代的 V-循环收敛性}
\label{sec:ch06-v-cycle-convergence}
\setcounter{thmcount}{0}
\setcounter{equation}{0}

本节考虑对称 V-循环算法. 换言之, 考虑第~\ref{sec:ch06-multigrid-algorithm} 节中 $p=1$ 且 $m_1=m_2=m$ 的算法.

我们的目标是证明此算法对任意 $m$ 都收敛, 因此不能再只作二重网格分析后接扰动论证. 必须直接分析初始误差 $z-z_0$ 与最终误差 $z-MG(k,z_0,g)$ 之间的关系.

\MBSourceItem{1}
\begin{Definition}
\label{def:ch06-v-cycle-error-operator}
误差算子 $E_k:V_k\longrightarrow V_k$ 递归定义为
\[
\MathBookFormulaID{formula-ch06-v-cycle-error-operator}
E_1=0,
\qquad
E_k=R_k^m[I-(I-E_{k-1})P_{k-1}]R_k^m,
\quad k>1.
\]
\end{Definition}

\MathBookSourcePage{178}
下一个引理说明 $E_k$ 把第 $k$ 层迭代的最终误差与初始误差联系起来.

\MBSourceItem{2}
\begin{Lemma}
\label{lem:ch06-error-operator-represents-final-error}
若 $z,g\in V_k$ 满足 $A_kz=g$, 则
\MBSourceItem{3}
\begin{equation}
\MathBookFormulaID{formula-ch06-error-operator-final-error}
\label{eq:ch06-error-operator-final-error}
E_k(z-z_0)=z-MG(k,z_0,g),
\qquad k\geq1.
\end{equation}
\end{Lemma}
\begin{Proof}
用归纳法证明. 当 $k=1$ 时, 由于 $MG(1,z_0,g)=A_1^{-1}g=z$, 结论显然. 假设式~\eqref{eq:ch06-error-operator-final-error} 对 $k-1$ 成立. 令 $q\in V_{k-1}$ 为粗网格上残差方程的精确解, 即 $A_{k-1}q=\bar g$. 因为 $q_1=MG(k-1,0,\bar g)$, 归纳假设给出 $q-q_1=E_{k-1}(q-0)$. 因而由引理~\ref{lem:ch06-coarse-residual-ritz-projection},
\MBSourceItem{4}
\begin{equation}
\MathBookFormulaID{formula-ch06-first-coarse-correction}
\label{eq:ch06-first-coarse-correction}
\begin{aligned}
q_1&=(I-E_{k-1})q\\
&=(I-E_{k-1})P_{k-1}(z-z_m).
\end{aligned}
\end{equation}
因此
\[
\MathBookFormulaID{formula-ch06-error-operator-induction-chain}
\begin{aligned}
z-MG(k,z_0,g)
&=z-z_{2m+1}\\
&=R_k^m(z-z_{m+1})\\
&=R_k^m(z-(z_m+q_1))\\
&=R_k^m(z-z_m-(I-E_{k-1})P_{k-1}(z-z_m))\\
&=R_k^m(I-(I-E_{k-1})P_{k-1})(z-z_m)\\
&=R_k^m(I-(I-E_{k-1})P_{k-1})R_k^m(z-z_0)\\
&=E_k(z-z_0).
\end{aligned}
\]
\end{Proof}

算子 $E_k$ 显然是线性的. 它关于 $a(\cdot,\cdot)$ 半正定则不那么明显. 为证明这一事实并简化收敛性证明, 先给出一个引理, 其证明留作习题(参见习题~\ref{exr:ch06-06}).

\MBSourceItem{5}
\begin{Lemma}
\label{lem:ch06-relaxation-energy-symmetry}
对所有 $v,w\in V_k$,
\[
\MathBookFormulaID{formula-ch06-relaxation-energy-symmetry}
a(R_kv,w)=a(v,R_kw).
\]
\end{Lemma}

\MBSourceItem{6}
\begin{Proposition}
\label{prop:ch06-error-operator-positive-semidefinite}
当 $k\geq1$ 时, $E_k$ 关于 $a(\cdot,\cdot)$ 对称半正定.
\end{Proposition}
\begin{Proof}
用归纳法证明. 当 $k=1$ 时显然成立. 假设命题对 $k-1$ 成立. 先证明 $E_k$ 关于 $a(\cdot,\cdot)$ 对称:

\MathBookSourcePage{179}
\[
\MathBookFormulaID{formula-ch06-error-operator-symmetry-chain}
\begin{aligned}
a(E_kv,w)
&=a(R_k^m(I-P_{k-1}+E_{k-1}P_{k-1})R_k^mv,w)
&&\text{由定义~\ref{def:ch06-v-cycle-error-operator}}\\
&=a((I-P_{k-1}+E_{k-1}P_{k-1})R_k^mv,R_k^mw)
&&\text{反复使用引理~\ref{lem:ch06-relaxation-energy-symmetry}}\\
&=a((I-P_{k-1})R_k^mv,R_k^mw)
 +a(E_{k-1}P_{k-1}R_k^mv,R_k^mw)\\
&=a(R_k^mv,(I-P_{k-1})R_k^mw)
 +a(E_{k-1}P_{k-1}R_k^mv,R_k^mw)\\
&=a(R_k^mv,(I-P_{k-1})R_k^mw)
 +a(E_{k-1}P_{k-1}R_k^mv,P_{k-1}R_k^mw)
&&\text{由定义~\ref{def:ch06-ritz-projection}}\\
&=a(R_k^mv,(I-P_{k-1})R_k^mw)
 +a(P_{k-1}R_k^mv,E_{k-1}P_{k-1}R_k^mw)
&&\text{由归纳假设}\\
&=a(R_k^mv,(I-P_{k-1})R_k^mw)
 +a(R_k^mv,E_{k-1}P_{k-1}R_k^mw)
&&\text{由定义~\ref{def:ch06-ritz-projection}}\\
&=a(R_k^mv,(I-P_{k-1}+E_{k-1}P_{k-1})R_k^mw)\\
&=a(v,R_k^m(I-P_{k-1}+E_{k-1}P_{k-1})R_k^mw)
&&\text{再次使用引理~\ref{lem:ch06-relaxation-energy-symmetry}}\\
&=a(v,E_kw).
\end{aligned}
\]
再证明 $E_k$ 关于 $a(\cdot,\cdot)$ 半正定:
\[
\MathBookFormulaID{formula-ch06-error-operator-positivity-chain}
\begin{aligned}
a(E_kv,v)
&=a(R_k^m(I-(I-E_{k-1})P_{k-1})R_k^mv,v)\\
&=a((I-(I-E_{k-1})P_{k-1})R_k^mv,R_k^mv)\\
&=a(R_k^mv,R_k^mv)
 -a((I-E_{k-1})P_{k-1}R_k^mv,R_k^mv)\\
&=a(R_k^mv,R_k^mv)
 -a((I-E_{k-1})P_{k-1}R_k^mv,P_{k-1}R_k^mv)\\
&=a(R_k^mv,R_k^mv)-a(P_{k-1}R_k^mv,P_{k-1}R_k^mv)\\
&\qquad+a(E_{k-1}P_{k-1}R_k^mv,P_{k-1}R_k^mv)\\
&=a((I-P_{k-1})R_k^mv,(I-P_{k-1})R_k^mv)\\
&\qquad+a(E_{k-1}P_{k-1}R_k^mv,P_{k-1}R_k^mv)\\
&\geq0,
\end{aligned}
\]
最后一步使用了归纳假设.
\end{Proof}

\MBSourceItem{7}
\begin{Lemma}[V-循环的光滑性质]
\label{lem:ch06-v-cycle-smoothing-property}
有
\[
\MathBookFormulaID{formula-ch06-v-cycle-smoothing-property}
a((I-R_k)R_k^{2m}v,v)
\leq\frac1{2m}a((I-R_k^{2m})v,v).
\]
\end{Lemma}
\begin{Proof}
\MathBookSourcePage{180}
先注意, 对 $0\leq j\leq l$,
\MBSourceItem{8}
\begin{equation}
\MathBookFormulaID{formula-ch06-relaxation-monotonicity}
\label{eq:ch06-relaxation-monotonicity}
\begin{aligned}
a((I-R_k)R_k^lv,v)
&=(A_k(I-R_k)R_k^lv,v)_k\\
&=(A_k\Lambda_k^{-1}A_kR_k^lv,v)_k\\
&=\frac1{\Lambda_k}(R_k^lA_kv,A_kv)_k\\
&\leq\frac1{\Lambda_k}(R_k^jA_kv,A_kv)_k
&&\text{由~\eqref{eq:ch06-relaxation-operator-properties}}\\
&=a((I-R_k)R_k^jv,v).
\end{aligned}
\end{equation}
因此
\[
\MathBookFormulaID{formula-ch06-v-cycle-smoothing-telescope}
\begin{aligned}
(2m)a((I-R_k)R_k^{2m}v,v)
&=a((I-R_k)R_k^{2m}v,v)+\cdots+a((I-R_k)R_k^{2m}v,v)\\
&\leq a((I-R_k)v,v)+a((I-R_k)R_kv,v)+\cdots\\
&\qquad+a((I-R_k)R_k^{2m-1}v,v)
&&\text{由~\eqref{eq:ch06-relaxation-monotonicity}}\\
&=a((I-R_k^{2m})v,v),
\end{aligned}
\]
其中最后一步为望远镜消去.
\end{Proof}

\MBSourceItem{9}
\begin{Proposition}
\label{prop:ch06-v-cycle-error-operator-bound}
令 $m$ 为光滑步骤数. 则
\MBSourceItem{10}
\begin{equation}
\MathBookFormulaID{formula-ch06-v-cycle-error-operator-bound}
\label{eq:ch06-v-cycle-error-operator-bound}
a(E_kv,v)\leq\frac{C^*}{m+C^*}a(v,v)
\qquad\forall v\in V_k,
\end{equation}
其中 $C^*$ 是与 $k$ 无关的正常数.
\end{Proposition}
\begin{Proof}
先估计 $a((I-P_{k-1})R_k^mv,(I-P_{k-1})R_k^mv)$:
\MBSourceItem{11}
\begin{equation}
\MathBookFormulaID{formula-ch06-projected-relaxed-error-bound}
\label{eq:ch06-projected-relaxed-error-bound}
\begin{aligned}
&a((I-P_{k-1})R_k^mv,(I-P_{k-1})R_k^mv)\\
&\quad=\|(I-P_{k-1})R_k^mv\|_E^2\\
&\quad\leq Ch_k^2\|R_k^mv\|_{2,k}^2
&&\text{由推论~\ref{cor:ch06-approximation-property}}\\
&\quad=Ch_k^2(A_k^2R_k^mv,R_k^mv)_k
&&\text{由~\eqref{eq:ch06-mesh-dependent-norm-scale}}\\
&\quad=Ch_k^2a(A_kR_k^mv,R_k^mv)
&&\text{由~\eqref{eq:ch06-discrete-operator-definition}}\\
&\quad=Ch_k^2\Lambda_ka((I-R_k)R_k^mv,R_k^mv)
&&\text{由定义~\ref{def:ch06-relaxation-operator}}\\
&\quad\leq Ca((I-R_k)R_k^mv,R_k^mv)
&&\text{由引理~\ref{lem:ch06-spectral-radius-bound}}\\
&\quad=Ca((I-R_k)R_k^{2m}v,v)
&&\text{由引理~\ref{lem:ch06-relaxation-energy-symmetry}}\\
&\quad\leq\frac Cm a((I-R_k^{2m})v,v)
&&\text{由引理~\ref{lem:ch06-v-cycle-smoothing-property}}.
\end{aligned}
\end{equation}
令 $C^*$ 为最后一个不等式中的常数 $C$. 下面对 $k\geq1$ 归纳证明式~\eqref{eq:ch06-v-cycle-error-operator-bound}. 当 $k=1$ 时, 因 $E_1=0$, 结论显然. 假设式~\eqref{eq:ch06-v-cycle-error-operator-bound} 对 $k-1$ 成立. 令 $\gamma=C^*/(m+C^*)$, 因而 $\gamma=(1-\gamma)C^*/m$. 则

\MathBookSourcePage{181}
\[
\MathBookFormulaID{formula-ch06-v-cycle-error-bound-induction}
\begin{aligned}
a(E_kv,v)
&=a(R_k^mv,R_k^mv)-a(P_{k-1}R_k^mv,P_{k-1}R_k^mv)\\
&\qquad+a(E_{k-1}P_{k-1}R_k^mv,P_{k-1}R_k^mv)\\
&=a((I-P_{k-1})R_k^mv,(I-P_{k-1})R_k^mv)\\
&\qquad+a(E_{k-1}P_{k-1}R_k^mv,P_{k-1}R_k^mv)\\
&\leq a((I-P_{k-1})R_k^mv,(I-P_{k-1})R_k^mv)\\
&\qquad+\gamma a(P_{k-1}R_k^mv,P_{k-1}R_k^mv)\\
&=(1-\gamma)a((I-P_{k-1})R_k^mv,(I-P_{k-1})R_k^mv)\\
&\qquad+\gamma a((I-P_{k-1})R_k^mv,(I-P_{k-1})R_k^mv)\\
&\qquad+\gamma a(P_{k-1}R_k^mv,P_{k-1}R_k^mv)\\
&=(1-\gamma)a((I-P_{k-1})R_k^mv,(I-P_{k-1})R_k^mv)\\
&\qquad+\gamma a(R_k^mv,R_k^mv)\\
&\leq(1-\gamma)C^*m^{-1}a((I-R_k^{2m})v,v)
 +\gamma a(R_k^mv,R_k^mv)\\
&=\gamma a((I-R_k^{2m})v,v)+\gamma a(R_k^mv,R_k^mv)\\
&=\gamma a(v,v).
\end{aligned}
\]
\end{Proof}

命题~\ref{prop:ch06-v-cycle-error-operator-bound} 立即蕴含 Braess \& Hackbusch 的如下经典结果.

\MBSourceItem{12}
\begin{Theorem}[V-循环第 $k$ 层迭代的收敛性]
\label{thm:ch06-v-cycle-level-convergence}
令 $m$ 为光滑步骤数. 则
\[
\MathBookFormulaID{formula-ch06-v-cycle-level-convergence}
\|z-MG(k,z_0,g)\|_E
\leq\frac{C^*}{m+C^*}\|z-z_0\|_E.
\]
因此, 对任意 $m$, 第 $k$ 层迭代都是收缩映射, 且收缩因子与 $k$ 无关.
\end{Theorem}
\begin{Proof}
由引理~\ref{lem:ch06-error-operator-represents-final-error}, 只需证明
\[
\MathBookFormulaID{formula-ch06-error-operator-energy-contraction}
\|E_kv\|_E\leq\frac{C^*}{m+C^*}\|v\|_E
\qquad\forall v\in V_k.
\]
由命题~\ref{prop:ch06-error-operator-positive-semidefinite}, 可对 $E_k$ 应用谱定理. 设 $0\leq\alpha_1\leq\cdots\leq\alpha_{n_k}$ 为 $E_k$ 的特征值, $\eta_1,\eta_2,\ldots,\eta_{n_k}$ 为相应特征向量, 且 $a(\eta_i,\eta_j)=\delta_{ij}$. 写成 $v=\sum\nu_i\eta_i$. 则 $a(E_kv,v)=\sum\alpha_i\nu_i^2$. 因而式~\eqref{eq:ch06-v-cycle-error-operator-bound} 蕴含
\[
\MathBookFormulaID{formula-ch06-error-operator-spectrum-bound}
0\leq\alpha_1\leq\cdots\leq\alpha_{n_k}
\leq\frac{C^*}{m+C^*}.
\]
所以
\MathBookSourcePage{182}
\[
\MathBookFormulaID{formula-ch06-error-operator-energy-proof}
\begin{aligned}
\|E_kv\|_E^2
&=a(E_kv,E_kv)\\
&=\sum_{i=1}^{n_k}\alpha_i^2\nu_i^2\\
&\leq\left(\frac{C^*}{m+C^*}\right)^2\sum_i\nu_i^2\\
&=\left(\frac{C^*}{m+C^*}\right)^2\|v\|_E^2.
\end{aligned}
\]
\end{Proof}
